\section{Лемма Лагранжа}

Задача: найти базис, в котором матрица Грама симметричной формы принимает простейший вид.
Оказывается, что симметричную форму всегда можно привести к диагональному виду.

\begin{lemma}[Лагранжа]
    Пусть $\mathrm{char}\, k \neq 2$, $B$ --- симметричная билинейная форма на $V$.
    Тогда существует базис, в котором матрица Грама $B$ диагональна.

    \begin{proof}
        Индукция по $\dim V$.

        \textbf{База.} Если $B \equiv 0$, то в любом базисе $G = 0$, и утверждение тривиально.

        \textbf{Переход.} Пусть $B \not\equiv 0$. Покажем, что $\exists\, e_1 \in V\colon B(e_1, e_1) \neq 0$.

        Если бы $B(v, v) = 0$ для всех $v \in V$, то для любых $u, v$:
        \[
            B(u + v,\, u + v) - B(u, u) - B(v, v) = 2B(u, v) = 0,
        \]
        и значит $B(u, v) = 0$ для всех $u, v$ (так как $\mathrm{char}\, k \neq 2$). Противоречие с $B \not\equiv 0$.

        Итак, $\exists\, e_1\colon B(e_1, e_1) \neq 0$.
        Возьмём $U = \langle e_1 \rangle$. Сужение $B\big|_U$ невырождено, поскольку $B(e_1, e_1) \neq 0$.

        По теореме об ортогональном дополнении $V = U \oplus U^\perp$.
        По индукционному предположению $U^\perp$ имеет базис $e_2, \dots, e_n$, в котором $B\big|_{U^\perp}$ имеет диагональную матрицу Грама.

        Тогда в базисе $e_1, e_2, \dots, e_n$ матрица Грама $B$ диагональна:
        на диагонали стоят $B(e_i, e_i)$, а $B(e_1, e_j) = 0$ при $j \geq 2$, так как $e_j \in U^\perp$.
    \end{proof}
\end{lemma}


\begin{cor}
    Пусть $\mathrm{char}\, k \neq 2$ и поле $k$ \emph{квадратично замкнуто} (для любого $\alpha \in k$ существует $\sqrt{\alpha} \in k$).
    Тогда для любой симметричной формы $B$ существует базис, в котором
    \[
        G = \mathrm{diag}(\underbrace{1, \dots, 1}_{r},\, \underbrace{0, \dots, 0}_{n - r}), \text{то есть } G = E \oplus O
    \]
    где $r = \mathrm{rk}\, B$.
    \begin{proof}
        По лемме Лагранжа найдём базис $e_1, \dots, e_n$ с диагональной матрицей Грама $G = \mathrm{diag}(d_1, \dots, d_n)$.
        Перенумеруем так, чтобы $d_1, \dots, d_r \neq 0$ и $d_{r+1} = \dots = d_n = 0$.

        Для $i \leq r$ положим $f_i \coloneqq e_i / \sqrt{d_i}$, для $i > r$ положим $f_i \coloneqq e_i$. Тогда:
        \[
            B(f_i, f_i) = \frac{B(e_i, e_i)}{d_i} = 1 \quad (i \leq r), \qquad B(f_i, f_i) = 0 \quad (i > r).
        \]
    \end{proof}
\end{cor}


\section{Теорема Дарбу}

Теперь рассмотрим каноническую форму кососимметричных (невырожденных) форм.
В отличие от симметричного случая, здесь не требуется никаких условий на поле.

\begin{definition}
    \emph{Симплектическое пространство} --- это пространство $(V, B)$ с невырожденной кососимметричной билинейной формой $B$,
    матрица Грама которой в некотором базисе имеет вид
    \[
        J = \begin{pmatrix} 0 & E \\ -E & 0 \end{pmatrix},
    \]
    где $E$ --- единичная матрица размера $n \times n$, и $\dim V = 2n$.
\end{definition}


\begin{theorem}[Дарбу]
    Пусть $V$ --- конечномерное пространство над произвольным полем, 
    $B$ --- невырожденная кососимметричная билинейная форма. Тогда $(V, B)$ изометрически изоморфно симплектическому пространству.

    В частности, $\dim V$ чётна.

    \begin{proof}
        Индукция по $\dim V$.

        \textbf{$\dim V = 1$ невозможно.} Если $\dim V = 1$ с базисом $e_1$,
        то $B(\alpha e_1, \beta e_1) = \alpha\beta\, B(e_1, e_1) = 0$ для всех $\alpha, \beta$ (кососимметричность).
        Это противоречит невырожденности.

        \textbf{База: $\dim V = 2$.} Так как $B$ невырождена, $\exists\, u, v\colon B(u, v) \neq 0$. Положим
        \[
            e_1 \coloneqq u, \qquad e_2 \coloneqq \frac{v}{B(u, v)}.
        \]
        Тогда $B(e_1, e_2) = 1$, $B(e_2, e_1) = -1$, и на диагонали нули (кососимметричность). Матрица Грама:
        \[
            G = \begin{pmatrix} 0 & 1 \\ -1 & 0 \end{pmatrix}.
        \]
        Заметим, что $u$ и $v$ линейно независимы: иначе $u = \alpha v$, и тогда $B(u, v) = \alpha B(v, v) = 0$ --- противоречие.

        \textbf{Переход.} Аналогично находим $e_1, e_2$ с $B(e_1, e_2) = 1$. Пусть $U = \langle e_1, e_2 \rangle$.

        Сужение $B\big|_U$ невырождено, так как
        \[
            \det G_U = \det \begin{pmatrix} 0 & 1 \\ -1 & 0 \end{pmatrix} = 1 \neq 0.
        \]
        По теореме об ортогональном дополнении $V = U \oplus U^\perp$. Сужение $B\big|_{U^\perp}$ остаётся невырожденным кососимметричным.

        По индукционному предположению $U^\perp$ имеет базис $e_3, e_4, \dots, e_{2n}$,
        в котором матрица Грама состоит из блоков $\bigl(\begin{smallmatrix}0 & 1 \\ -1 & 0\end{smallmatrix}\bigr)$.

        В базисе $e_1, e_2, \dots, e_{2n}$ матрица Грама всего $V$ блочно-диагональна:
        \[
            G = \mathrm{diag}\left(
                \begin{pmatrix} 0 & 1 \\ -1 & 0 \end{pmatrix},
                \begin{pmatrix} 0 & 1 \\ -1 & 0 \end{pmatrix},
                \dots
            \right).
        \]

        Наконец, перенумеруем базис: $f_1 = e_1,\, f_2 = e_3,\,
        \dots,\, f_n = e_{2n-1}$ и $f_{n+1} = e_2,\, f_{n+2} = e_4,\, \dots,\, f_{2n} = e_{2n}$. В этом базисе матрица Грама принимает вид
        \[
            J = \begin{pmatrix} 0 & E \\ -E & 0 \end{pmatrix}.
        \]
    \end{proof}
\end{theorem}


\begin{remark}
    В отличие от симметричного случая (лемма Лагранжа), теорема Дарбу не требует никаких ограничений на поле:
    ни характеристику, ни квадратичную замкнутость. Каноническая форма невырожденной кососимметричной формы одна и та же над любым полем.
\end{remark}
